\section{An asymptotic obstruction to constant ramp tracking}
\label{sec:revisit-ramp-asymptotic}
This section proves an asymptotic obstruction for the actual nonlinear
graph trajectory. The proof first identifies a linear amplification
mechanism, then justifies its graph limit using the cooling energy and
level symmetry. It is not needed by any positive algorithm theorem above.

Fix $\theta=1/64$, $\alpha=\theta^2$, $\chi=1-\theta$,
$\eta=1-\theta/2$, $\beta=\chi/(1+\theta)$, and
$c=(1-\alpha)/2$. In degree-density coordinates define
\[
 Q_B=\begin{pmatrix}(1+\alpha)/2&-c\\-c/2&(1+\alpha)/2\end{pmatrix},
 \qquad C_B=I-Q_B=c\begin{pmatrix}1&1\\1/2&1\end{pmatrix}.
\]
These matrices are self-adjoint in the inner product with weights $(1,2)$.
They arise from an even/odd pair of levels of $\mathcal T(b,2,L)$ as
$b\to\infty$. The coupling from the next pair's first density into this
pair's second density is $c/2$ in $I-Q$; the reverse coupling vanishes
in that limit. This limiting calculation does not supply a quotient oracle
to a local algorithm.

\begin{lemma}[Exact frequency amplification in a level pair]
\label{lem:revisit-pair-gain}
For a positive pair obeying the linear recurrence, divide its PPR density
error by $r_k$ and take temporal differences. At frequency $\zeta$ on
the unit circle, set $z=\eta\zeta$ and
\[
 A(z)=c(1+\beta-\beta/z),\qquad
 H(z)=\frac{A(z)^2}{2(z-A(z))^2-A(z)^2}.
\]
The transfer from the next pair's first error difference to the current
pair's first error difference is $H(z)$. At $\zeta=(12+5\mathrm i)/13$,
\[
 |H(\eta\zeta)|^2
 =\frac{51067952908040250000}{5105426406330152281}>9.
\]
All poles of the normalized causal transfer lie strictly inside the unit
circle.
\end{lemma}
\begin{proof}
For error $e_k=t-x_k$ on a region where the kinetic floor is inactive,
\[
 e_{k+1}=(I-Q)((1+\beta)e_k-\beta e_{k-1})+\alpha r_k\one.
\]
After division by $r_{k+1}=\eta r_k$, temporal differences remove the
constant forcing. The pair's frequency matrix is
\[
 \begin{pmatrix}z-A&-A\\-A/2&z-A\end{pmatrix},
\]
with right-hand side $(0,A/2)^\trans$ times the next pair's first
coordinate. Solving its two equations gives the transfer formula.
The displayed strict inequality is rational arithmetic; it is independently
checked by \texttt{front\_limit\_probe.py}.

The eigenvalues of $C_B$ are $c(1\pm1/\sqrt2)$. The unnormalized
second-order poles have moduli
$\chi\sqrt{(1\pm1/\sqrt2)/2}$, so their largest modulus is
$\sigma=\chi\cos(\pi/8)<\chi<\eta$. Division by $\eta$ proves stability.
\end{proof}

\paragraph{The isolated activation profile.}
Use $Q_B$, source $(1,0)^\trans$, and regularizer $\lambda_\tau=\eta^\tau$
in the orthant recurrence. Extend it by zero for integer $\tau\leq0$;
that extension is consistent because both raw entries are nonpositive
when $\lambda_\tau\geq1$. Put
\[
 t_B=Q_B^{-1}(1,0)^\trans,\qquad
 U(\tau)=\alpha\eta^{-\tau}(t_B-x(\tau)),\qquad
 d(\tau)=U_1(\tau)-U_1(\tau-1).
\]

\begin{lemma}[The activation signal excites an amplifying frequency band]
\label{lem:revisit-front-signal}
The sequence $d$ is nonzero and decays exponentially in both time
directions. If $f_j$ is the result of passing it through $j$ copies of
the transfer in \cref{lem:revisit-pair-gain}, then there is a constant
$a>0$, independent of $j$, such that
\[
 \sup_{\tau\in\mathbb Z}|f_j(\tau)|\geq a\,3^j/\sqrt j
 \qquad(j\geq1).
\]
\end{lemma}
\begin{proof}
The pair has a Stieltjes Hessian with spectrum in $[\alpha,1]$ in its
weighted metric, and $Q_B\one\geq\alpha\one$.
Apply the principal-operator auxiliary energy argument of
\cref{sec:revisit-pruning}. Here the weighted source mass is one and
$\lambda_0=1$, so its bank recursion is
$B^+\leq\chi B+(\lambda-\lambda^+)$ and $B\leq\lambda$.
Consequently both primal and kinetic states tend to the strictly positive
$t_B$. From some finite time onward the kinetic floor is inactive.

Thereafter the profile is a linear combination of the poles of modulus
at most $\sigma$ and the particular response $\eta^\tau v$, where
\[
 v=(\eta I-(1+\beta-\beta/\eta)C_B)^{-1}\one>0.
\]
Thus $U(\tau)\to\alpha v$, and its temporal difference decays at rate
at most a polynomial times $(\sigma/\eta)^\tau$. In the negative time
direction, $U=\alpha\eta^{-\tau}t_B$, which also decays. Moreover
$\sum_\tau d(\tau)=\alpha v_1>0$, so $d$ is nonzero.

Its bilateral transform $D(\zeta)=\sum_\tau d(\tau)\zeta^{-\tau}$
is analytic in an annulus containing the unit circle. The transfer is
analytic there after shrinking the annulus if necessary. By
\cref{lem:revisit-pair-gain}, an open arc has $|H(\eta\zeta)|>3$.
The nonzero analytic $D$ cannot vanish on the entire arc. A closed
subarc therefore has $|D|\geq a_0>0$ and $|H|\geq3$.
Parseval's identity gives $\|f_j\|_2\geq a_1 3^j$.

To turn this into a maximum bound, choose circles of radii
$r<1<R$ inside the common annulus. Cauchy's coefficient estimates give
\[
 |f_j(\tau)|\leq C B_r^j r^\tau\quad(\tau\geq0),\qquad
 |f_j(-\tau)|\leq C B_R^j R^{-\tau}\quad(\tau\geq0),
\]
for fixed finite $C,B_r,B_R$. Choose a constant $L$ sufficiently large
that both tails $|\tau|>Lj$ have squared sum less than
$a_1^2 3^{2j}/2$ for all sufficiently large $j$.
The remaining $O(j)$ coefficients therefore contain a value of magnitude
at least $a_2 3^j/\sqrt j$. Decreasing the constant covers the finitely
many smaller $j$, whose output is nonzero.
\end{proof}

\subsection{Justifying the separated-front limit on finite unit trees}
Fix an integer $j\geq1$ and let
\[
 \gamma=\frac{c^2}{2((1+\alpha)/2)^2-c^2},\qquad
 b_K=\operatorname{nearestInteger}(\gamma\eta^{-K}),
 \qquad G_{j,K}=\mathcal T(b_K,2,2j+4).
\]
Run the actual point-source ramp on $G_{j,K}$, and write $E_{j,K}(\tau)$
for its root PPR density error divided by $r_{jK+\tau}$.
The graph-limit statement is
\begin{equation}
 E_{j,K}(\tau)-E_{j,K}(\tau-1)\longrightarrow f_j(\tau)
 \quad\text{as }K\to\infty,
 \quad\text{for every fixed }j,\tau.
 \tag{FL}\label{eq:revisit-front-limit}
\end{equation}
We prove this limit below. All limits first fix $j$ and any finite time
window, and then send $K$ to infinity. Constants may depend on these fixed
quantities and on the fixed $\alpha$, but not on $K$ or $b_K$.

Write $T_\ell$ for the exact PPR degree density on level $\ell$, and
$X_{\ell,k},Z_{\ell,k}$ for the actual primal and kinetic densities.
Every level is an automorphism orbit fixing the seed, so these densities
are constant within a level. Let $W_\ell$ be that level's original degree
volume. For pair $i=(2i,2i+1)$, $0\leq i\leq j$, put
\[
 h_i=\frac\alpha b\left(\frac\gamma b\right)^i,
 \qquad \mathbf T_i=(T_{2i},T_{2i+1})^\trans.
\]

\begin{lemma}[Static scales and coordinatewise energy control]
\label{lem:revisit-level-scales}
For fixed $j$ and $0\leq i\leq j$, as $b\to\infty$,
\[
 \mathbf T_i/h_i\longrightarrow t_B,\qquad
 W_{2i}/b^{i+1}\longrightarrow1,\qquad
 W_{2i+1}/b^{i+1}=2.
\]
At every actual ramp step, on every level,
\begin{equation}
 |X_{\ell,k}-T_\ell|,
 |Z_{\ell,k}-T_\ell|\leq2\sqrt{r_k/W_\ell}.
 \label{eq:revisit-level-energy}
\end{equation}
\end{lemma}
\begin{proof}
Here $\mu=\alpha$, so \cref{eq:revisit-response-mass} and
$Q\succeq\alpha I$ imply both
$\|x_k-t\|_2^2\leq4r_k$ and $\|z_k-t\|_2^2\leq4r_k$.
Grouping the nonnegative squared terms by level proves
\cref{eq:revisit-level-energy}. The level volumes follow by counting:
level $2i$ has $b^i$ vertices of degree $b+1$ except that the root has
degree $b$; level $2i+1$ has $b^{i+1}$ vertices of degree two.

For completeness, the static asymptotics follow by elimination from the
leaves. Set $R_\ell=T_\ell/T_{\ell-1}$ and $q_0=(1+\alpha)/2$.
At the last leaf, $R=c/q_0<1$. At other odd and even levels,
respectively, elimination gives
\[
 R_\ell=\frac{c/2}{q_0-(c/2)R_{\ell+1}},\qquad
 R_\ell=\frac{c/(b+1)}{q_0-(cb/(b+1))R_{\ell+1}}.
\]
Backward induction gives $0<R_\ell<1$. Every nonterminal even ratio
is $O(1/b)$, because its denominator is at least $\alpha$.
Thus for $i\leq j$, the next even level is nonterminal and
$R_{2i+1}=c/(2q_0)+O(1/b)$. It follows that
$R_{2i}R_{2i-1}=\gamma/b+O(1/b^2)$ for $1\leq i\leq j$.
Finally $T_0=(\alpha/b)/(q_0-cR_1)$, yielding the stated $t_B$ scale.
The extra pair before the terminal leaf ensures these formulas hold
through pair $j$.
\end{proof}

\begin{lemma}[Activation separation without an assumed history]
\label{lem:revisit-activation-separation}
For each $i\leq j$ there is a fixed integer $C_i$ such that, for all
sufficiently large $K$, both kinetic coordinates in pair $i$ are strictly
positive at every ramp time $k\geq iK+C_i$.
For every fixed integer $M$, all levels beyond pair $j$ remain identically
zero up to time $jK+M$, for sufficiently large $K$.
Moreover pair $j$ and all later levels are zero up to time $jK-1$.
\end{lemma}
\begin{proof}
Since $b_K\eta^K\to\gamma$, one has
\[
 r_{iK+C}=(1/b_K)\eta^{iK+C},\qquad
 r_{iK+C}W_{2i}\longrightarrow\gamma^i\eta^C,
 \qquad T_{2i}W_{2i}\longrightarrow\alpha\gamma^i(t_B)_1,
\]
and analogous positive limits on the odd coordinate. Choose $C_i$ large
enough that $2\sqrt{r_kW_\ell}<T_\ell W_\ell$ at $k=iK+C_i$
on both levels. This remains true at later ramp times because $r_k$
decreases. Equation~\eqref{eq:revisit-level-energy} proves strict
kinetic positivity, and hence inactivity of the kinetic projection floor.

The inactivity claims require an argument from the start of the run.
Suppose all levels from an even level $2i$ onward are zero at the incoming
state. Only its first level can receive a nonzero raw forcing, namely
$cY_{2i-1,k}/(b+1)-\alpha r_k$ before division by $\theta$.
Since $Y$ is a convex combination of $X$ and $Z$,
\[
 \frac{cY_{2i-1,k}}{(b+1)\alpha r_k}
 \leq\frac{cT_{2i-1}}{(b+1)\alpha r_k}
   +\frac{2c}{(b+1)\alpha\sqrt{r_kW_{2i-1}}}.
\]
Both terms increase as $r_k$ decreases, so it suffices to check the last
time under consideration. At $i=j$, $k=jK-1$, the first term tends to
$\eta<1$ and the second is $O(b^{-1/2})$.
At $i=j+1$, $k=jK+M$, both terms are $O(1/b)$.
For sufficiently large $b$, the raw forcing is therefore nonpositive
at every earlier time. Induction from the zero initialization proves
the claims for these separating levels and all levels after them.
\end{proof}

\begin{lemma}[The actual current front has the isolated profile limit]
\label{lem:revisit-current-front}
For every fixed integer $\tau$,
\[
 \frac{(X_{2j,jK+\tau},X_{2j+1,jK+\tau})^\trans}{h_j}
       \longrightarrow x(\tau),\qquad
 \frac{(Z_{2j,jK+\tau},Z_{2j+1,jK+\tau})^\trans}{h_j}
       \longrightarrow z(\tau).
\]
Consequently this pair's normalized PPR error tends to $U(\tau)$.
\end{lemma}
\begin{proof}
The pair is zero through $\tau=-1$, and all later pairs are zero on any
fixed time window, by \cref{lem:revisit-activation-separation}.
Its incoming source is $cY_{2j-1,k}/(b+1)$. The static part divided by
$h_j$ tends to one. The error in that source, divided by $h_j$, is
$O(b^{-1/2})$ on each fixed window by
\cref{eq:revisit-level-energy}. Its internal matrix tends to $Q_B$, and
$\alpha r_{jK+\tau}/h_j\to\eta^\tau$.
Induction over the finitely many steps from $\tau=-1$ now applies the
continuity of all operations, including the coordinatewise maximum.
Negative fixed times are already zero. Finally
$h_j/r_{jK+\tau}\to\alpha\eta^{-\tau}$ and
\cref{lem:revisit-level-scales} identify the limiting PPR error.
\end{proof}

\begin{theorem}[The graph-limit statement holds]
\label{thm:revisit-front-limit}
Equation~\eqref{eq:revisit-front-limit} holds for the actual nonlinear
ramp on the finite canonical unit graphs $G_{j,K}$.
\end{theorem}
\begin{proof}
Write $\mathbf E_{i,k}$ for pair $i$'s density error divided by $r_k$.
Fix a last time $T=jK+M$. First we establish a bound on every
$\mathbf E_{i,k}$, $i\leq j$, $0\leq k\leq T$, that is uniform in $K$.
For pair $j$ this follows directly from
\cref{eq:revisit-level-energy}, because $r_TW_\ell$ has a positive
limit. For each $i<j$, the same estimate bounds its error uniformly up
to $t_i=iK+C_i$. After $t_i$, both kinetic coordinates are positive
by \cref{lem:revisit-activation-separation}, so its error recurrence
is exactly linear.

The within-pair matrix in $I-Q$ is $C_B$ at the root and
\[
 C_{i,b}=c\begin{pmatrix}1&b/(b+1)\\1/2&1\end{pmatrix}
 \quad(i>0).
\]
Its normalized second-order recurrence has poles uniformly bounded in
modulus by $\sigma/\eta<1$. For example, its eigenvalues lie between
$c(1-1/\sqrt2)$ and $c(1+1/\sqrt2)$ and the quadratic temporal poles
have modulus $\sqrt{\beta\lambda}/\eta$. Fix
$\sigma/\eta<s_0<1$. The corresponding impulse kernels and homogeneous
propagators have norms at most $C s_0^m$, uniformly for sufficiently
large $b$. This follows by the resolvent formula on $|z|=s_0$ for the
four-dimensional state matrix, whose entries converge and whose spectrum
is strictly inside that circle.

The next pair enters with a bounded coefficient $c/2$, while the previous
pair enters with coefficient $c/(b+1)$. Summing these uniformly
absolutely summable kernels, including the one-step delayed inputs, gives
\[
 M_i\leq A_i+L M_{i+1}+(L/b)M_{i-1},\qquad
 M_i:=\max_{0\leq k\leq T}\|\mathbf E_{i,k}\|_\infty,
 \quad i<j,
\]
where $M_{-1}=0$, $M_j$ is already bounded, and all $A_i,L$ are
independent of $K$. The constants include the uniformly bounded states
at $t_i,t_i-1$ and the constant regularization forcing.
The nonnegative matrix of this finite system tends to a strictly upper
triangular matrix. For sufficiently large $b$ its spectral radius is
less than one, and its inverse $\sum_{m\geq0}N_b^m$ is nonnegative
and uniformly bounded. Multiplication by this inverse proves the claimed
uniform bounds on all $M_i$. This step prevents a hidden growth factor
from the earlier fronts.

Now shift time by $jK$. For $i<j$ the time since $t_i$ tends to infinity,
so the uniformly bounded homogeneous initial contribution tends to zero.
The previous-pair contribution is $O(1/b)$ by the bounds just proved.
The remaining causal kernels converge to those of $C_B$ and have a common
exponentially summable majorant. Starting with the current-front limit in
\cref{lem:revisit-current-front}, dominated convergence of the causal
sums proves the older-pair limits successively for $i=j-1,\ldots,0$.
To apply this on an arbitrarily long past, first truncate the kernel at
a fixed length; its remaining tail is uniformly small by the same
$M_i$ bounds, and then let that length tend to infinity.

The limiting recurrence is exactly the positive-pair recurrence used in
\cref{lem:revisit-pair-gain}, with constant regularization forcing.
Temporal differences remove that forcing. Each pair therefore contributes
one copy of the transfer $H$ to the current front's difference $d$.
At the root the result is $f_j$, proving
\cref{eq:revisit-front-limit}.
\end{proof}

\begin{corollary}[No graph-uniform constant ramp-tracking bound]
\label{cor:revisit-unbounded-tracking}
No graph-uniform constant $C$ can bound the ramp's normalized PPR error
by $Cr_k$, even at the fixed
$\alpha=1/4096$ and with target $\rho=1/(1024\vol(G))$.
\end{corollary}
\begin{proof}
Choose $j$ so that \cref{lem:revisit-front-signal} exceeds $2C$ at some
finite integer $\tau$. The limit then makes one of the two neighboring
values of $|E_{j,K}|$ exceed $C$ for sufficiently large $K$.
At these times $r=\Theta(b_K^{-j-1})$, whereas
$\rho=\Theta(b_K^{-j-2})$, so the ramp has not reached its target.
\end{proof}

\paragraph{Measured checks and proof scope.}
The limit-profile audit agrees at 80 and 100 decimal digits. For four
old pairs, the maximum discrepancy on $0\leq\tau\leq800$ between the
finite-tree profile and its limit decreases from about $2.03\cdot10^{-5}$
at $b=6{,}481{,}113$ to $3.82\cdot10^{-8}$ at $b=3{,}440{,}713{,}565$.
These are diagnostics for (FL), not its proof. The proof above uses the
global energy only to establish levelwise sign separation and initial
bounds, then obtains a uniform normalized-error bound through the finite
stable input system. It does not assume monotonicity, lower order, or
nonnegative residuals along the ramp. The obstruction is specific to
omitting the final accuracy phase of this recurrence; it is not a lower
bound against all local algorithms. The earlier unresolved route and its
resolution are preserved in \path{RAMP_ASYMPTOTICS.md}.
